\documentclass{beamer}

\usepackage{amsmath, amssymb, mathtools}
\newcommand\norm[1]{\left\lVert#1\right\rVert}
\usepackage{multirow, array, cellspace}
    \setlength{\cellspacetoplimit}{2pt}
    \setlength{\cellspacebottomlimit}{2pt}

\newcommand{\appropto}{\mathrel{\vcenter{
  \offinterlineskip\halign{\hfil$##$\cr
    \propto\cr\noalign{\kern2pt}\sim\cr\noalign{\kern-2pt}}}}}
    
\usepackage{graphicx}
\usepackage[style=authoryear]{biblatex}
\renewcommand*{\nameyeardelim}{\addcomma\addspace}

\addbibresource{../references.bib}

\usepackage[dvipsnames]{xcolor}
\definecolor{first}{HTML}{0077BB}
\definecolor{second}{HTML}{009988}
\definecolor{third}{HTML}{CC3311}
\newcommand\firstplace[1]{\textcolor{first}{\textbf{#1}}}
\newcommand\secondplace[1]{\textcolor{second}{\underline{\smash{}#1}}}
\newcommand\thirdplace[1]{\textcolor{third}{\textit{#1}}}

% Import the UNCMath theme.
% UNC is the default color theme.
\usepackage{source_files_slides/beamerthemeUNCMath}

\logo{
    
\includegraphics[height=3em]{STOR_712_Project/source_files_slides/oldwell_cmyk-eps-converted-to.pdf}
}

% A macro for making the title frame.
% Removes the bottom bar and logo temporarily.
% If you don't want these in other frames, 
% you could try mimicking this.
\newcommand{\titleframe}{
{
\setbeamertemplate{headline}{}
\setbeamertemplate{footline}{}
\begin{frame}
\titlepage
\end{frame}
}
}

%%%%%%%%%%%%%%%%%%%%%%%%%%%%%%%%%%%%%%%%%%
%
% Information for the title frame.
%
\title{
    Can We Improve \textbf{Personalized Federated Learning} With \textbf{Learning Rate Scheduling}?
}
\subtitle{
    \bigskip
    Learning from Models rather than Data
}
\author{
    Joe Lavond 
    and Peter Lin
}
\date{\today}
%%%%%%%%%%%%%%%%%%%%%%%%%%%%%%%%%%%%%%%%%%
%
% Start of document.
%
\begin{document}

% This is the titleframe, do it this way 
% so that you don't have the logo on the title page.
\titleframe


\begin{frame}<0>{Introduction to Federated Learning}
    \begin{definition}[Federated Learning (FL)]
        An iterative procedure allowing \textcolor{blue}{multiple users} to build a shared robust machine learning model \textcolor{red}{without sharing data} (e.g. healthcare). 
    \end{definition}

    Key pieces:
    \begin{itemize}
        \item Data: Multiple users with their data
        \item Training: Users each have a local copy of the model for training
        \item \textbf{Communication}: Send model information rather than data \\
            \quad (privacy and resource demands) \\
    \end{itemize}

\end{frame}


\begin{frame}{Review of FL}
    FL is an iterative process of training rounds. See Figure~\ref{fig: fl}.
    
    \begin{figure}
        \centering
        \includegraphics[width = .8\textwidth]{STOR_712_Project/images/fedavg.png}
        \caption{Single Round of Federated Training}
        \label{fig: fl}
    \end{figure}
    %
    The most common aggregation scheme for Step 4 in Figure~\ref{fig: fl} is FedAvg \parencite{fedavg}, which averages returning models.
\end{frame}


\begin{frame}<0>{Personalized Federated Learning (pFL)}

    \textcolor{red}{Q}: Is a single model appropriate if users have very different data sets? \\
    
    \smallskip
    \textcolor{blue}{A}: \textcite{pfl-bench} argues for individualized models tailored to each user's data.
    
    \bigskip
    How to personalize?
    \begin{itemize}
        \item Main works Ditto \parencite{ditto}, FedBN \parencite{fedbn}, and FedEM \parencite{fedem} require additional computation and memory
        \item Simplest solution is fine-tuning from the single Federated Learning model
    \end{itemize}
        
    Is it possible to find a single model that, after fine-tuning, performs similarly to the more complex solutions?
    
\end{frame}


\begin{frame}{User Local Training}

    Within a round of FL, each user performs $K$ updates to the global model
    \begin{align}
        \theta_{k + 1} = \theta_{k} - \textcolor{red}{\alpha} \nabla F(\theta_{k}) 
    \end{align}

    Motivated by Taylor Expansion of \textcite{reptile}, we propose learning rate scheduling (within each training round)
    \begin{align}
        & \phi_{k + 1} = \phi_{k} - \textcolor{blue}{\alpha \beta_k} \nabla F(\phi_{k}) \label{eq: lr-decay} 
        %\\ & \phi_K - \phi_0 \propto \sum_{k = 1}^K \beta_{k - 1} \nabla F(\phi_{k - 1}) 
    \end{align}

\end{frame}


\begin{frame}[allowframebreaks]{Taylor Analysis}

    Expansion of our proposed update where the gradients and Hessians are evaluated using the current global model.
    \begin{align}
        & \mathbb{E} \left[ \phi_K - \phi_0 \right] \appropto - B(K) \textcolor{blue}{\mathbb{E} \left[ \Tilde{g}_j \right]} + \left( \alpha \sum_{j = 1}^K \beta_{(j - 1)} B(j - 1) \right) \textcolor{blue}{\mathbb{E} [ \Tilde{H}_j \Tilde{g}_h ]}
    \end{align}

    \begin{enumerate}
        \item $\mathbb{E} \left[ \Tilde{g}_j \right]$ - minimize joint training objective (initial model success)
        \item $\mathbb{E} [ \Tilde{H}_j \Tilde{g}_h ]$ - maximizes inner product between gradients of different mini-batches (rapid personalization), see Equation~\ref{eq: AvgGradInner}
    \end{enumerate}
    
    \begin{align}
        & \mathbb{E} [ \Tilde{H}_k \Tilde{g}_l ] 
            = \frac{1}{2} \mathbb{E} [ \Tilde{H}_k \Tilde{g}_l + \Tilde{H}_l \Tilde{g}_k ] 
            = \frac{1}{2} \mathbb{E} \left[ \frac{\partial}{\partial \phi_0} \left( \Tilde{g}_k^T \Tilde{g}_l \right) \right] \label{eq: AvgGradInner}
    \end{align}


% \framebreak
%     Using exponential local decay $\beta_j = \beta^j$, then 

%     {
    
%     \begin{table}[h]
%         \centering
        
%         \begin{tabular}{|Sl|c|c|c|c|c|}
%         \hline
%             \multirow{2}{*}{Method} & 
%                 FedSGD & 
%                 \textcolor{blue}{FedDecay} & 
%                 FedAvg & 
%                 FOMAML \\
%             & 
%                 $\beta = 0$ & 
%                 \textcolor{blue}{$\beta \in \mathopen(0, 1\mathclose)$} & 
%                 $\beta = 1$ & 
%                 $\beta_j = I(j = K)$ \\
%             \hline
%             Ratio $\dfrac{\#( \mathbb{E}[ \Tilde{H}_j \Tilde{g}_h ] )}{\#( \mathbb{E} [ \Tilde{g}_j ])}$ & 
%                 $\dfrac{0}{1}$ &
%                 $\dfrac{\left( \beta - \beta^{K} \right) \eta}{1 - \beta^2}$ & 
%                 $\dfrac{(K - 1)\eta}{2}$ &
%                 $(K - 1)\eta$ \\
%             Limit $(K \rightarrow \infty)$ & 
%                 0 & 
%                 $\dfrac{\beta \eta}{1 - \beta^2}$ &
%                 $\infty$ &
%                 $\infty$ \\
%             \hline
%         \end{tabular}
        
%         \caption{\small
%             Ratios of counts for $\mathbb{E} [ \Tilde{H}_j \Tilde{g}_h ]$ to $\mathbb{E} [ \Tilde{g}_j ]$ for several methods.
%         }
%     \end{table}
%     }
    
% \end{frame}

\framebreak
    Using exponential local decay $\beta_j = \beta^j$, then 
    \medskip
    \resizebox*{\textwidth}{!}{       
        \begin{tabular}{|Sl|c|c|c|c|c|}
        \hline
            \multirow{2}{*}{Method} & 
                FedSGD & 
                \textcolor{blue}{FedDecay} & 
                FedAvg & 
                FOMAML \\
            & 
                $\beta = 0$ & 
                \textcolor{blue}{$\beta \in \mathopen(0, 1\mathclose)$} & 
                $\beta = 1$ & 
                \\
            \hline
            Coef Ratio & 
                $\dfrac{0}{1}$ &
                $\dfrac{\left( \beta - \beta^{K} \right) \eta}{1 - \beta^2}$ & 
                $\dfrac{(K - 1)\eta}{2}$ &
                $(K - 1)\eta$ \\
            Limit $(K \rightarrow \infty)$ & 
                0 & 
                $\dfrac{\beta \eta}{1 - \beta^2}$ &
                $\infty$ &
                $\infty$ \\
            \hline
        \end{tabular}
    }
    The ratios of coefficients for $\mathbb{E} [ \Tilde{H}_j \Tilde{g}_h ]$ and $\mathbb{E} [ \Tilde{g}_j ]$ represent the trade-off between emphasis on how easily it can be fine-tuned for each user vs. initial model success.
    
\end{frame}



\begin{frame}[allowframebreaks]{Convergence Results}

    \textbf{Assumptions}:

    \medskip
    For all users $i \in \{1, \ldots, M\}$, the local objective function $F_i$ 
    \begin{enumerate}
        \item is $\mu_i$-strongly convex 
        \item has $L_i$-Lipschitz continuous gradients.
    \end{enumerate}

    \medskip
    Furthermore, their stochastic gradients
    \begin{enumerate}
        \item have bounded variance
        \item have uniformly bounded expected squared norm, $\mathbb{E} \norm{g_k}^2 \leq M$
    \end{enumerate}


\framebreak
    Need learning rate decay over rounds for convergence.
    \begin{align}
        & \eta_t = \alpha_{\lfloor t / K \rfloor} \beta_{(t \bmod K)} 
            \textrm{ where } \alpha_j = \frac{c}{j + d} \textrm{ for some } c, d > 0 
    \end{align}

    \begin{theorem}
        \label{thm: full}
        Under previous assumptions, even with non-iid local datasets, then
        \begin{align*}
            & \mathbb{E} \left[ F ( \theta_g^{(N)} ) \right] - F^* \leq \dfrac{\kappa}{\textcolor{blue}{N} + d - 2} \times \left( C_0 + C_1 \mathbb{E} \norm{\theta_g^{(1)} - \theta_g^*}^2 \right) 
        \end{align*}
        for $N$ communication rounds with $K$ local update steps.
    \end{theorem}
    
\end{frame}


\begin{frame}[allowframebreaks]{Numerical Experiments}

%     Q: How does our learning rate scheduling impact model performance?
%     \begin{itemize}
%         \item Hold out data and users to measure generalization
%         \item Evaluate across an image, text, and graph benchmark data set
%         \item Compare with fixed learning rate FedAvg and personalized methods.
%     \end{itemize}

%     \medskip
%     Notation: 
%     \begin{itemize}
%         \item $\Bar{Acc}$, $\Breve{Acc}$, and $\sigma_{Acc}$ represent the equally-weighted average, bottom ten percentile, and standard deviation of test set accuracy.
%         \item Denote metrics for new users who did not participate in federated training with an additional tick suffix ($^\prime$).
%     \end{itemize}

%\framebreak

    Metrics: $\Bar{Acc}$, $\Breve{Acc}$, and $\sigma_{Acc}$ represent the equally-weighted average, bottom tenth percentile, and standard deviation of test set accuracy
    
    \begin{table}[htb]
        \resizebox{\textwidth}{!}{
            \begin{tabular}{|l|rrr|rrr|rrr|}
            \hline
                & 
                    \multicolumn{3}{c|}{FEMNIST (Image)} & 
                    \multicolumn{3}{c|}{SST (Text)} & 
                    \multicolumn{3}{c|}{PUBMED (Graph)} \\ 
                Method & 
                    $\Bar{Acc}$ & $\Breve{Acc}$ & $\sigma_{Acc}$ &
                    $\Bar{Acc}$ & $\Breve{Acc}$ & $\sigma_{Acc}$ &
                    $\Bar{Acc}$ & $\Breve{Acc}$ & $\sigma_{Acc}$ \\
            \hline
                Ditto & 
                    0.5672 & 0.4444 & 0.0921 &
                    0.4746 & 0.0000 & 0.4010 &
                    0.2442 & - & - \\
                FedBN & 
                    0.9059 & 0.8302 & 0.0544 &
                    \textbf{0.8030} & \textbf{0.6667} & 0.1275 &
                    \textbf{0.8004} & - & - \\
                FedEM & 
                    \textbf{0.9175} & 0.8333 & 0.0526 &
                    0.7404 & 0.6379 & 0.2141 &
                    0.7879 & - & - \\
                pFedMe & 
                    0.9036 & \textbf{0.8438} & 0.0751 &
                    0.7785 & 0.6406 & 0.1179 &
                    0.7932 & - & - \\
            \hline
                FOMAML & 
                    0.8989 & 0.8113 & 0.0604 &
                    0.7680 & 0.6667 & 0.1057 &
                    0.7950 & - & - \\
                FedAvg & 
                    0.9055 & \textbf{0.8491} & 0.0530 &
                    0.7680 & 0.6667 & 0.1057 &
                    0.7914 & - & - \\
                \textcolor{blue}{\textbf{FedDecay}} & 
                    \textbf{0.9152} & 0.8421 & 0.0485 &
                    % 0.9194 & 0.8421 & 0.0476 &
                    \textbf{0.8101} & \textbf{0.6724} & 0.1087 &
                    % 0.7708 & 0.6667 & 0.0987 & 
                    \textbf{0.8039} & - & - \\
            \hline
            \end{tabular}
        }
        \caption{Generalization to new users who did not participate in federated training after fine-tuning}
        \label{tab: metrics--unseen}
    \end{table}

\end{frame}


\begin{frame}[allowframebreaks]{References}
    \printbibliography
\end{frame}



\end{document}